\section{Budgeting, compaction and task evidence delivery}
\subsection{Exact provider counting before inference}
\texttt{countContext()} requests \texttt{/v1/responses/input\_tokens}. At 90\% of the configured capacity, \texttt{compactContext()} calls the selected model's \texttt{/compact} endpoint. Core retains the \emph{entire} returned canonical window, including opaque compaction/reasoning items, checks the 16 MiB serialization bound and recounts. If the new window remains above threshold or compaction is unsupported, inference fails visibly while saved history/evidence remains available.

\texttt{contextBudget()} permits an operator override of 8,192--2,000,000 tokens. Code contains capacity entries for three named models and a 128,000-token fallback; these are implementation configuration, not independently established provider guarantees. Output is bounded by the smaller of 8,192 tokens and 8\% of capacity. Listing a model does not prove that counting, compaction or thinking is compatible.
\subsection{Incremental task context}
\texttt{taskContextDelta()} fingerprints operation status/result/error/attempts. Initial context includes request, plan/criteria and up to 20 conversation messages, each limited to 4,000 characters. Subsequent turns receive changed observations and changed tool metadata plus the latest 40-operation ledger, ten completed-step labels and deterministic check statuses. Results above 12,000 characters become marked previews. Updates target a 64,000-character changed-evidence batch; a single first oversized change can exceed that target, so it is not a universal hard cap.

Delivery fingerprints and the private window commit together under owner fencing. \texttt{readTaskEvidence()} reads only this task's stored operation IDs and applies JSON-pointer/offset/limit paging. Arrays are limited to 50 items and reduced by character size; an oversized single item becomes a preview and needs a more specific pointer. Paging reads stored facts; it never re-executes a remote operation or counts as new external proof.
\subsection{Long-task checkpoints and pollution control}
At a per-prompt turn/token/time boundary, \texttt{summarizeTaskContext()} requests a factual memo from the same model (2,048 output tokens, 16,000-character returned-text limit). \texttt{continuePrompt()} commits the memo, resets prompt baselines, retains lifetime accounting and requeues. Failed summaries do not overwrite saved context. Immutable criteria, grant identities and write ledger survive independently of prose.

Instructions label server descriptions and retrieved content untrusted and preserve hypotheses, IDs and declined/uncertain operations. This mitigates prompt pollution but does not eliminate injection. Stale observations can survive a summary; industrial freshness must be checked explicitly. There is no semantic chat-memory index or learned personal-memory subsystem.
\verdict{Bounded deltas and stored-result paging reduce repeated context growth while preserving exact evidence access. Summaries remain lossy, and the current implementation does not deterministically enforce observation freshness.}
\source{\path{core/src/tasks/task-context.ts}: \texttt{taskContextDelta()}, \texttt{readTaskEvidence()}; \path{core/src/llm/llm.service.ts}: \texttt{summarizeTaskContext()}; \path{core/src/llm/context-window.ts}.}
